\section{问题重述}

通用神经网络处理器在单核内部需要同时协调计算节点、数据搬运节点以及多级片上缓存。给定一个带有前驱依赖的计算图，每个操作节点具有固定的执行流水、执行周期和所使用的 buffer；每个 buffer 又由 \texttt{ALLOC}/\texttt{FREE} 节点界定逻辑生命周期，并被指定到 L1、UB、L0A、L0B 或 L0C 等缓存类型。不同合法调度顺序会改变 buffer 生命周期的重叠程度、物理地址复用方式以及各执行流水的并行关系，因此内存与时间性能彼此耦合。

本文将三问统一理解为对同一计算图逐层增加约束的优化问题。

\begin{enumerate}[label=\textbf{问题\arabic*:},leftmargin=5.5em]
    \item 在保持 DAG 前驱关系、buffer 生命周期与 L0 硬约束合法的条件下重新排列节点，最小化 L1+UB 的峰值逻辑驻留量。
    \item 在合法拓扑调度基础上加入真实缓存容量与连续地址约束，为所有 buffer 分配物理地址；当空间不足或碎片阻塞时允许插入 SPILL，并最小化由此增加的 DDR 数据搬运量。
    \item 在问题二得到的物理内存与 SPILL 方案基础上，进一步考虑 CUBE、VECTOR、MTE1、MTE2、MTE3 和 FIXP 等执行流水的并行与串行关系。通过地址重排和调度优化减少总执行周期，并分析数据搬运量与执行周期之间的可选权衡。
\end{enumerate}

六组测试数据分别覆盖 Matmul、FlashAttention 和 Conv 三类计算图，规模从千级节点到三万余节点不等。因此算法必须同时满足两个要求：一方面严格遵守题面给出的内存、SPILL 和时序语义；另一方面避免依赖对全排列空间的穷举，而应使用可以在大规模稀疏 DAG 上稳定运行的结构化启发式与局部搜索。